% =====================================================================
% DRAFT APPENDIX: Symmetries, Representations, and Invariance
% Rewritten (NOT lifted) from Bronstein et al. (2021) GDL proto-book, sec. 3
% \citep{Bronstein2021_geometricDL_book}, per CKJ margin comment p.16.
% Depth: EXPANSIVE on 3.1, SHORT on 3.2 & 3.3, MEDIUM on 3.4; 3.5 (the
% blueprint) lives in the main text. House style: light on abstract math,
% heavy on physical-system analogies. Group-theory statements refer to
% \citep{Bronstein2021_geometricDL_book} and/or Zee's "Group Theory in a
% Nutshell for Physicists" \citep{Zee2016_grouptheory_nutshell}; other
% citations (Klein's Erlangen Programme, Mallat for wavelets) replicated from
% Bronstein and can be promoted to bib entries if we keep them.
% TODO (CKJ): revisit this appendix later -- run-in headings (svmult \runinhead
% vs orphaned \paragraph), the subsection title duplicating the section title,
% and general structure/typography.
% =====================================================================

\section{Appendix: Symmetries, Representations, and Invariance}
\label{app:symmetries}

This appendix collects the geometric vocabulary the main text leans on. We follow the framing of \citet{Bronstein2021_geometricDL_book}, but keep the mathematics light and the physics close; readers wanting the full group-theory machinery will find a physicist-friendly treatment in \citet{Zee2016_grouptheory_nutshell}.

\subsection{Symmetries, representations, and invariance}
\label{app:sym:basics}

A \emph{symmetry} is a transformation that leaves some property of a system unchanged. Physicists meet symmetries before they meet the word: the outcome of an experiment does not depend on \emph{where} the bench sits (translation), \emph{which way} it faces (rotation), or \emph{when} it is run (time translation); a molecule's energy does not care how the molecule is oriented; a set of identical particles does not care in what order we list them. Each ``does not care'' is a symmetry, and -- as Noether taught us -- each continuous one is shadowed by a conserved quantity.

Symmetries compose: doing one and then another is again a symmetry, every symmetry can be undone, and ``do nothing'' is a symmetry. Those three facts are exactly the axioms of a \emph{group}, the natural home for a system's symmetries. We keep the definition formal, since everything else hangs off it:

\begin{definition}[Group]
A \emph{group} is a set $G$ with a composition rule
$\circ : G\times G \to G$ (written $g\circ h = gh$) obeying:
\begin{itemize}
  \item \textbf{Associativity:} $(gh)k = g(hk)$ for all $g,h,k\in G$;
  \item \textbf{Identity:} a unique $e\in G$ with $eg = ge = g$ for all $g$;
  \item \textbf{Inverse:} each $g$ has a unique $g^{-1}$ with $gg^{-1}=g^{-1}g=e$;
  \item \textbf{Closure:} $gh\in G$ for all $g,h\in G$.
\end{itemize}
\end{definition}

Commutativity is \emph{not} required: in general $gh \neq hg$ (rotating then translating differs from translating then rotating). Groups for which $gh=hg$ always hold are called \emph{Abelian}. Large groups are often built from a handful of \emph{generators} -- e.g.\ the symmetries of an equilateral triangle are generated by one rotation and one reflection, and that same six-element group is also the group of permutations of three labelled corners. This little coincidence is the seed of a recurring theme: the \emph{same} abstract group can act on very different objects.

\paragraph{Group actions and representations.} We rarely care about a group in the abstract; we care about how it \emph{acts} on our data. If the group acts on the underlying domain $\Omega$ (rotating points in space, relabelling nodes of a graph), it automatically acts on signals defined over that domain (rotating an image, permuting node features). The actions that matter for us are \emph{linear}, and a linear action is precisely a \emph{representation}: a rule $\rho$ that assigns to each group element $g$ an invertible matrix $\rho(g)$ such that $\rho(gh)=\rho(g)\rho(h)$ -- composing transformations becomes multiplying matrices. The physicist's reflex is the right one here: this is exactly how a scalar, a vector, or a higher tensor ``transforms under rotation'' (the rotation matrices, or for higher spins the Wigner $D$-matrices), and the dimension of $\rho$ is just the size of the feature it acts on, not of the group itself. (\citealp{Zee2016_grouptheory_nutshell} develops representations from this physical standpoint.)

\paragraph{Invariance vs.\ equivariance.} Imposing a symmetry on a learned function $f$ is a powerful inductive bias: it shrinks the space of candidate functions to those that already respect the symmetry, so less has to be learned from data. Two cases recur throughout the review:

\begin{definition}[Invariance and equivariance]
A function $f$ is \emph{$G$-invariant} if $f(\rho(g)x) = f(x)$ for all $g\in G$ -- the output ignores the transformation. It is \emph{$G$-equivariant} if $f(\rho(g)x) = \rho(g)\,f(x)$ -- the output transforms in lockstep with the
input.
\end{definition}

The physical reading is immediate. A molecule's total energy is rotation- and permutation-\emph{invariant} (rotate the molecule, relabel its atoms -- the energy is unmoved); the force on each atom is rotation-\emph{equivariant} (rotate the molecule and every force vector rotates with it) and permutation-\emph{equivariant} (relabel the atoms and the list of forces is relabelled the same way). The same split organizes graph tasks: a graph-level prediction should be permutation-invariant, while node- and edge-level predictions should be permutation-equivariant. Convolutional layers are the canonical equivariant building block -- shifting the input shifts the output feature map by the same amount -- and a final global pooling turns equivariance into invariance. That ``stack equivariant layers, then pool'' pattern is the geometric-deep-learning blueprint discussed in the main text.

\subsection{Levels of structure: isomorphisms and automorphisms}
\label{app:sym:iso}

What counts as a symmetry depends on \emph{how much structure} we insist on preserving. A bare set only has a size, so its symmetries are all relabellings (bijections). Demand that nearness be preserved and you keep only continuous maps; demand smoothness, only diffeomorphisms; demand distances, only \emph{isometries} (rotations, translations, reflections). Each extra requirement carves out a \emph{subgroup} of the previous group: more structure preserved means fewer allowed transformations. This nesting is exactly Klein's Erlangen-Programme view of geometry (projective $\supset$ affine $\supset$ Euclidean), and it is the lens we use to decide which symmetries a given astrophysical problem actually has.

One distinction is worth stating plainly because it is routinely blurred. An \emph{automorphism} maps an object to \emph{itself} while preserving its structure -- this is what we have been calling a symmetry. An \emph{isomorphism} maps between \emph{two different} objects that are structurally the same. For graphs this is the everyday statement that two graphs are ``the same graph drawn differently'': a graph isomorphism is a node relabelling that preserves every connection (and non-connection), while a graph \emph{automorphism} is such a relabelling of a graph onto itself -- a genuine internal symmetry. This is the precise sense in which a GNN's permutation symmetry says ``the labels are arbitrary; only the relations matter.''

\subsection{Deformation stability: when symmetries are only approximate}
\label{app:sym:deformation}

Exact symmetries are an idealization. The real world is noisier: in a video, several objects each drift their own way, so no single global shift relates consecutive frames; a galaxy image may be mildly warped by lensing or seeing yet still show the same galaxy. What we really want is not exact invariance to a clean group, but \emph{approximate} invariance to transformations that are \emph{close} to that group. The fix is to measure how far a transformation $\tau$ sits from the symmetry group $G$ with a cost $c(\tau)$ (zero on $G$ itself) and to ask only for \emph{deformation stability},
%
\begin{equation}
\| f(\rho(\tau)x) - f(x)\| \;\le\; C\,c(\tau)\,\|x\|,
\label{eq:def-stability}
\end{equation}
%
i.e.\ small deformations may change the output, but only a little. (Small deformations cannot simply be added to the group, because composing many small ones makes a large one -- they do not close.) This is the formal version of a demand we make repeatedly: a good model should degrade gracefully, not catastrophically, when the symmetry it assumes holds only approximately, and it connects to the \emph{soft}-constraint idea (imposing a symmetry as a prior rather than a hard architectural law).

The same idea applies when it is the \emph{domain}, not the signal, that deforms -- the case closest to our subject. A social or follow graph changes its edges over time; the cosmic web is a different graph at every redshift; a 3D shape flexes. One quantifies these by a distance between domains (e.g.\ graph edit distance, which vanishes for isomorphic graphs, or the Gromov--Hausdorff distance for manifolds) and asks that $f$ change slowly in that distance.

\subsection{Scale separation}
\label{app:sym:scale}

Deformation stability tames \emph{local} wobble, but on a large domain there are still far too many stable functions to learn efficiently -- the curse of dimensionality bites. The escape is \emph{scale separation}: most physical tasks are organized hierarchically, so we can build global understanding by composing local, multi-scale pieces.

The classical tool for ``which scales are present'' is the Fourier transform, which rewrites a signal as a sum of oscillations indexed by frequency, and which turns convolution into simple multiplication (the convolution theorem). Fourier descriptors capture \emph{global} properties (smoothness, overall power) and are perfectly matched to a clean global symmetry like translation -- the modulus of the Fourier transform, for instance, is exactly shift-invariant. But Fourier representations are \emph{fragile}: an arbitrarily small, high-frequency deformation can change them by an order-one amount, because each basis function is spread over the whole domain.

\emph{Wavelets} (and multi-scale representations generally) repair this by using filters localized in \emph{both} space and scale: fine details are captured by small-support atoms, coarse trends by large-support ones. Crucially, a multi-scale decomposition is approximately \emph{equivariant} to deformations (its error grows like the deformation size, $O(\epsilon)$, rather than $O(1)$), turning a globally unstable description into a family of locally stable features. Those features are not yet invariant; they must be combined progressively from fine to coarse -- which is precisely why useful architectures are \emph{deep} and \emph{hierarchical}, interleaving local operators with coarsening. The astrophysical resonance is hard to miss: structure formation, turbulence, and the cosmic web are all multi-scale, and a model that respects scale separation is matching the physics of its inputs. (See \citealp{Bronstein2021_geometricDL_book} and references therein, e.g.\ Mallat on wavelets, for the technical development.)

The graph inherits every piece of this story through the \emph{graph Laplacian} $L = D - A$, with $D$ the diagonal matrix of node degrees. $L$ acts as the discrete second derivative on the graph: apply it to a signal, and it returns zero wherever the signal runs smooth across the edges and large values wherever the signal jumps. Its eigenvectors are the graph's Fourier modes and its eigenvalues their frequencies, which cashes out the promise of Figure~\ref{fig:graph_anatomy} that the spectrum of $A$, or of its Laplacian, reveals communities, diffusion, and scale. Spectral graph convolutions filter in exactly this eigenbasis \citep{Bruna2013_spectral_networks, Defferrard2016_chebnet}, and the GCN of \citet{Kipf2016_GCN_semisupervised} descends from them as a fast local approximation. Message passing reaches the same operations without ever computing an eigendecomposition, which is why the main text keeps the spatial view.

% NOTE: 3.5 (the GDL blueprint: equivariant layers -> local pooling -> global
% pooling -> readout) is covered in the MAIN TEXT, not here (per CKJ).
